\section{A Siegel-unit period and lattice calculation}\label{sec:siegel}

This section constructs the integral smoothed classes used in the positive
comparison and determines their periods.  The three ingredients have distinct
roles.  A quotient that kills cusp differences gives a full integral
modular-symbol lattice.  Siegel-unit distributions give integral classes and
trace maps.  Explicit reciprocity and an archimedean calculation determine
their rational differential coordinates, including every smoothing and Euler
factor.  The integral divisibility of the normalized determinant will be
proved in the subsequent sections.

\subsection{The elliptic quotient and its relative lattice}

Fix an elliptic modular isogeny class of conductor \(N\), with normalized
newform \(f\).  Let \(E_0\) be the elliptic optimal quotient of \(J_1(N)\), so
that the quotient map has connected kernel.  Let \(C_{\mathrm{cusp}}\) be the
subgroup of \(E_0(\overline{\mathbb Q})\) generated by the images of all
geometric cusp differences, and put
\[
                         E'=E_0/C_{\mathrm{cusp}}.
\]
Manin--Drinfeld makes \(C_{\mathrm{cusp}}\) finite, and its Galois stability
makes this a \(\mathbb Q\)-isogeny \cite{ManinDrinfeld}.  If \(D_N\) is the cusp
set, the Abel map followed by this quotient gives
\[
                  \pi:(X_1(N),D_N)\longrightarrow(E',O)
\]
over \(\mathbb Q\).  Changing the base cusp changes the Abel map by a cusp
difference, so the resulting map is independent of that choice.

Allow \(A=1\) or \(A=2^k\), put \(L=NA\), and let
\[
                  \Pi:(X_1(L),D_L)\longrightarrow(E',O)
\]
be the \(A\)-degeneracy followed by \(\pi\).  The degeneracy quotients by the
cyclic subgroup of order \(A\) in the marking and pushes the marked point
forward.  At the mark \(1/L\), it is \(\tau\mapsto A\tau\).  Thus, for a
nonzero rational invariant differential \(\omega\) on \(E'\),
\[
 \Pi^*\omega=c_F F\,\frac{dq}{q},
 \qquad F(\tau)=A f(A\tau),\qquad q=e^{2\pi i\tau},
 \qquad c_F\in\mathbb R^\times.
\]
The form \(f\) has rational coefficients and trivial character and is pulled
back from \(X_0(N)\).  The marked levels are fine: weight-two cuspidality on
\(X_0(N)\) in particular gives \(N\geq5\).

Write
\[
 \Lambda_B=H_1(E'(\mathbb C),\mathbb Z),\qquad
 T=H^1(\overline E',\mathbb Z_2)(1).
\]
The principal pairing identifies
\(T=T_2(E')^\vee(1)\) with \(T_2E'\).  Choose a generator \(a\) of the
real-invariant sublattice of \(\Lambda_B\), oriented so that
\(\int_a\omega=\Omega_0>0\).  This is the connected real period.  The cycle
\(a\) is primitive and, under the comparison identifications, gives a basis
of \(T^{G_{\mathbb R}}\).  For \(\xi\in\mathrm{SL}_2(\mathbb Z)\), let
\(\delta_\xi\in H^1(E'(\mathbb C),\mathbb Z)\) be Poincaré dual to
\(\Pi_*\xi\{0,\infty\}\), with orientation fixed by
\[
       \int_{E'(\mathbb C)}\omega\wedge\delta_\xi
           =\int_{\Pi_*\xi\{0,\infty\}}\omega.
\]

\begin{lemma}[The integral relative lattice]\label{lem:siegel-lattice}
There is an integral Galois-equivariant map
\[
 H^1(\overline{Y_1(L)},\mathbb Z_2)(1)\longrightarrow T
\]
dual to pullback on relative cohomology.  When \(A=1\), there are finitely
many integers \(n_\xi\) such that
\[
              \delta=\sum_\xi n_\xi\delta_\xi,\qquad \delta(a)=1.
\]
\end{lemma}

\begin{proof}
Pullback by the map of pairs \(\Pi\) is integral.  Relative Poincaré duality
gives the displayed map; no rational Hecke projector is involved.  Its
differential coordinate for the dual exponential, with the Tate-twist
de Rham comparison, is paired by the cup invariant with the ordinary finite
Kummer logarithm on \(E'\).

For the lattice assertion, write \(E_0(\mathbb C)=\mathbb C/\Lambda_0\).
Connectedness of the kernel of \(J_1(N)\to E_0\) makes the map on integral
first homology onto \(\Lambda_0\).  The quotient by \(C_{\mathrm{cusp}}\)
replaces \(\Lambda_0\) by the lattice \(\Lambda_B\) generated by
\(\Lambda_0\) and lifts of the cusp images.  Relative cusp paths supply
exactly these additional generators.  Hence
\[
 H_1(X_1(N)(\mathbb C),D_N;\mathbb Z)\longrightarrow\Lambda_B
\]
is surjective.  The Farey triangulation generates the relative group by
Manin edges \(\xi\{0,\infty\}\).  Their images therefore generate
\(\Lambda_B\) integrally.

The invariant sublattice is the kernel of complex conjugation minus one
on a free integral lattice and is saturated; its generator \(a\) is
primitive.  Unimodularity of the intersection pairing on \(\Lambda_B\)
then makes evaluation on \(a\) surjective on \(H^1(E'(\mathbb C),\mathbb Z)\).
An integral combination of the Poincaré duals of the Manin edges consequently
has evaluation one.
\end{proof}

\subsection{Level patterns and integral smoothed classes}

We use the following two patterns of orders.  The dyadic valuation of \(M\)
will always be sufficiently large; for example, \(v_2(M)\geq6\) suffices for
the reciprocity calculation below.

\begin{enumerate}[label=(\roman*)]
\item In \emph{square order}, \(n_1=n_2=M\), with \(L\mid M\).  Set
\(l=e_2\xi\), where \(e_2=(0,1)\).
\item In \emph{rectangular order}, \(n_1=M\), \(n_2=bM\), where
\(b=N_{\mathrm{odd}}\), \((b,M)=1\), and \(2^{v_2(L)}\mid M\).  Set
\(l=(Ns,1)=e_2\xi\).  When filling the first row, also allow
\(n_1=b_1M\), where \(b_1\mid b\) is a product of some of the full prime
powers dividing \(b\).
\end{enumerate}

On \(Y(n_1,n_2)\), let \(P,Q\) be the independent marked points of the
indicated orders on the universal elliptic curve.  Independence means that
they define an embedding of the indicated product of cyclic groups.
The mark and determinant used for projection are
\[
 C=l_1(n_1/L)P+l_2(n_2/L)Q,\qquad
 \eta_M=e_M\bigl((n_1/M)P,(n_2/M)Q\bigr).
\]
Both coefficients in \(C\) are integers, including in the rectangular
case, and \(C\) has order \(L\).  The determinant \(\eta_M\) is a primitive
\(M\)-th root of unity.  Fix compatible primitive roots in algebraic,
dyadic, and complex embeddings, using \(e^{2\pi i/m}\) in the complex
embedding.

Trace to
\(Y_1(L)\times\operatorname{Spec}\mathbb Q(\mu_M)\), and then use
Lemma~\ref{lem:siegel-lattice}.  For a finite character
\(\nu\) of \(\operatorname{Gal}(\mathbb Q(\mu_M)/\mathbb Q)\), the summand
whose root is \(\zeta_M^\sigma\) receives weight \(\nu(\sigma)\).
This is an unnormalized trace: equivalently, take the group-coefficient
trace and send its basis vector at \(\sigma\) to \(\nu(\sigma)\).
For period calculations, trivialize the finite twist after restriction
using the specified embeddings.  Reversing the Galois convention replaces
the argument character consistently by its inverse.

Choose \(c,d>1\) prime to \(6n_2\).  On the variable elliptic curve \(E\),
write
\[
                {}_c g_P=P^*{}_c\theta_E,\qquad
                {}_d g_Q=Q^*{}_d\theta_E
\]
for the canonical smoothed Siegel units.  The divisor of
\({}_c\theta_E\) is \(c^2(0)-E[c]\); its norm under \([s]\) is itself
when \((s,c)=1\).  These are the normalized distributions of
\cite{Kato04}.  Norm compatibility under an isogeny of degree prime to
\(c\) follows from the divisor and uniqueness of this normalization:
the norms under \([2]\) and \([3]\) commute with the isogeny.  The Tate
formula in the proof below also checks the normalization directly.

At stage \(j\), replace the orders by \(2^j n_1,2^j n_2\) and let \(P',Q'\)
map to \(P,Q\) under multiplication by \(2^j\).  Before the mark and
determinant projection, the class is the edge image of the unnormalized
trace, at precision \(2^j\), of
\[
  ({}_c g_{P'})\cup({}_d g_{Q'})\ \otimes\
           e_{2^j}(n_1P',n_2Q')^\vee .                   \tag{K1}\label{eq:K1}
\]
Parentheses denote Kummer classes.  The dual basis sends the displayed
primitive root to \(1\), so the total coefficient is
\(\mathbb Z/2^j(1)\).

\begin{proposition}[Integral classes and dyadic compatibility]
\label{prop:siegel-integral}
The classes obtained from \eqref{eq:K1} are integral for support containing
the primes of \(2n_2\) and infinity.  They are compatible with the dyadic
division traces and, after the mark and determinant projection, with
corestriction as the dyadic part of \(M\) increases with \(c,d,A,\xi\)
fixed.  The projected smoothed classes on \(\mathbb Q\), with group or
character coefficients, admit integral lifts to ordinary positive
cohomology.
\end{proposition}

\begin{proof}
A further dyadic division gives a Cartesian product of the two sets of
division lifts.  Independence remains automatic, since the old tuple is
already a basis at 2.  At the old precision the inverse-root twist agrees.
Base change and the projection formula identify the trace of the cup on
this product with the cup of the two traces.  Siegel-unit distribution
then recovers the old class.

The evaluated units and their inverses are integral away from the level.
Horizontal valuations follow from their divisors.  For vertical
valuations, use smooth full-level models with cusps, on which every good
geometric fiber component meets a cusp.  The Tate expansion has leading
factors consisting of signs, roots of unity, and powers of \(1-\zeta\)
for nontrivial level roots; all are units away from the level.
A quotient of level can be checked after pullback to such a model.

The smooth proper pair and its residue sequence show that geometric
cohomology of the relative open curve is lisse away from this support
and vanishes above degree one.  The arithmetic degree-two to degree-one
edge map therefore lands integrally in
\[
 H^1\!\left(G_{\mathbb Q,S},
       H^1(\overline{Y(n_1,n_2)},\mathbb Z/2^j)(1)\right).
\]
Here degree-one cohomology of a lisse sheaf on the integer base is
fundamental-group cohomology; over a field this is the
Hochschild--Serre edge map.  Edge maps, maps of pairs, and traces commute.
Passage to compact degree-one classes uses finiteness, or the
Mittag--Leffler property in degree zero.  The lattice push and the
Shapiro coefficient push are integral, so neither requires a
Hecke-idempotent or group-projector denominator.

The same Cartesian-lift calculation proves compatibility when the base
dyadic part of \(M\) increases: the lower mark and determinant root agree.
For the positive lift, first retain the Shapiro coefficient at the full
cyclotomic level.  This field is imaginary, and the induced coefficient
has zero ordinary real \(H^1\).  After pushing coefficients, the global
class consequently has zero real restriction and lifts to the mapping
fiber defining positive cohomology.  Compatible lifts in a tower can be
chosen by compactness of the finite-degree cohomology groups at finite
precision.  Rational wedge with the real boundary basis is independent
of the choice of lift.
\end{proof}

\subsection{Full-level reciprocity and cusp calibration}

For a nonzero torsion label \(P=[x\tau+y]\), let \(E_P\) be the algebraic
weight-one Eisenstein form represented analytically by
\[
 -\frac{1}{2\pi i}\left.
 \sum_{(m,n)\in(x,y)+\mathbb Z^2}
 \frac{(\operatorname{Im}\tau)^w}
 {(m\tau+n)|m\tau+n|^{2w}}\right|_{w=0},
\]
where evaluation at zero means regular continuation.  Fix Fourier and
orientation conventions consistently throughout.

\begin{proposition}[The full-level differential]\label{prop:siegel-reciprocity}
In square order, the dual exponential of the inverse-root-twisted trace
\eqref{eq:K1}, before cuspidal projection, is on each determinant component
a scalar multiple of
\[
  (c^2 E_P-c E_{cP})(d^2 E_Q-d E_{dQ})\,\frac{dq}{q}.     \tag{K2}\label{eq:K2}
\]
The scalar is
\[
                              \pm M^{-2},               \tag{K3}\label{eq:K3}
\]
with one common sign determined by the cup, pairing, and orientation
conventions.  The statement holds for all \(c,d\) prime to \(6M\).
\end{proposition}

\begin{proof}
\emph{The explicit reciprocity input.}
Use the full-level product formula of \cite[Theorem~9.5]{Kato04},
specialized to \(k=2\) and \(r=r'=1\).  The containment of prime supports
required there is automatic at square level.  The condition
\(1\leq r\leq k-1\), with one of \(r,r'\) equal to \(k-1\), is satisfied.
We are in its case \(p\mid M\), so no additional Euler operator occurs.
The theorem applies at \(p=2\).  This is the full-level theorem;
the projected theorem \cite[Theorem~9.7]{Kato04} is not needed to obtain
\eqref{eq:K2}.

Initially take \(c,d\equiv1\pmod M\).  The symmetric-power coefficient
has degree zero.  The étale two-unit symbol regulator is the Kummer cup,
and the negative Tate specialization over the full dyadic division
tuples uses
\[
                         e_{2^j}(MP',MQ')^\vee.
\]
It is the determinant-to-the-power-\(-1\) moment.  Thus the test is a
full-level basis coset: every lift remains a basis tuple at 2.  On the
differential side no raising or polynomial-coefficient operator remains.
At weight one the shifted and symplectic Fourier-dual Eisenstein forms
agree under the consistent conventions; see
\cite[Proposition~3.11(1) and \S4.2]{Kato04}, or the Kronecker functional
equation.  Both labels are nonzero.  The smoothing on the first row is
\(c^2-c\,\mathrm{mult}_c^*\), and similarly on the second row.  For the
initial smoothing choices the two brackets in \eqref{eq:K2} are scalar
multiples of the individual forms.  Only this proportionality is needed
from explicit reciprocity; we now determine its scalar.

For completeness, the matrix form of this specialization uses a residue
coset
\[
                \begin{pmatrix}a&b\\ d'&e\end{pmatrix}\pmod M
\]
with primitive determinant.  Its two arguments are
\[
                  \frac{a\tau+b}{M},\qquad
                  \frac{d'\tau+e}{M}.
\]
At precision \(2^j\), the matrices lift this coset modulo \(2^jM\).
In the compatible cyclotomic basis, the inverse twist in
\eqref{eq:K1} is the inverse determinant of the lift modulo \(2^j\),
up to the fixed orientation.  The determinant need not be restricted
to a single lift.  This agrees with the distribution formula
\cite[\S2.3.3, Theorem~1.4, and Proposition~5.5]{Wang12} at
\((k,j)=(2,1)\); its depth condition \(v_p(M)\geq v_p(2p)\)
is satisfied here.  The Bloch--Kato dual exponential identification
used in the present proof is the original one in Kato's theorem.

\emph{The normalized theta function.}
Put
\[
 \Theta(t,q)=(1-t)\prod_{h\geq1}(1-q^ht)(1-q^h/t),
 \qquad \Theta(qt,q)=-t^{-1}\Theta(t,q).
\]
The normalized function is
\[
 {}_c\theta_E(t)=(-1)^{(c-1)/2}q^{(c^2-1)/12}
 t^{-c(c-1)/2}\frac{\Theta(t,q)^{c^2}}{\Theta(t^c,q)}.
\]
It is periodic and has the required divisor.  To check its norm
normalization, take the \([s]\)-preimages
\(t^{1/s}q^{h/s}\zeta_s^j\), \(0\leq h,j<s\), with \((s,c)=1\).
The numerator theta factors multiply, before taking the \(c^2\)-th
power, to \(\Theta(t,q)\).  After taking the product over \(j\),
the denominator factors are
\(\Theta(q^{ch}t^c,q^s)\), \(0\leq h<s\).  Write
\[
                    ch=h'+sk_h,\qquad 0\leq h'<s.
\]
The inverse quasiperiod factors have \(t^c\)-exponent and sign exponent
\[
                         \sum_h k_h=(c-1)(s-1)/2,
\]
and their \(q\)-exponent is
\[
 \sum_h\left(h'k_h+\frac{s k_h(k_h-1)}2\right)
 =\frac{(c^2-1)(s-1)(2s-1)}{12}
     -\frac{(c-1)s(s-1)}4.
\]
Together with the prefactors, this proves exact norm invariance.
The divisor and these normalizing norms determine the function
uniquely: their quotient is constant, and the norms for \([2]\) and
\([3]\) force that constant to have both its third and eighth powers
equal to one.

For \(0<x<1\), \(cx\notin\mathbb Z\), set
\[
 J_c(x)=c\lfloor cx\rfloor-\frac{c(c-1)}2
       =c^2B_1(x)-cB_1(\{cx\}),
 \qquad B_1(x)=x-\tfrac12.
\]
At \(t=q^x\zeta\), the leading coefficient is
\(\pm\zeta^{J_c(x)}\), and the \(q\)-order is
\[
 xJ_c(x)+\frac{c^2-1}{12}
            -\frac{\lfloor cx\rfloor(\lfloor cx\rfloor+1)}2.
\]
The final two terms are integers.

\emph{Residues and field transfer.}
On any determinant component choose the cusp represented in Tate
notation by
\[
                     P=q^{1/M},\qquad
                     Q=q^{1/M}\zeta_M^u,
                     \qquad u\in(\mathbb Z/M)^\times.
\]
Extend scalars to \(K_0=\mathbb Q_2(\mu_M)\).  Its width is \(M\).
Put \(M_j=2^jM\) and \(K_j=K_0(\mu_{M_j})\).  For each
\[
                a_1,a_2\in[0,M_j)\cap(1+M\mathbb Z),
\]
there is one cusp orbit over \(K_0\) with representative
\[
             P'=q^{a_1/M_j},\qquad
             Q'=q^{a_2/M_j}\zeta_{M_j}^u,
\]
width \(M_j\), and residue field \(K_j\).  Indeed, once the old labels
are fixed, translations by the old width identify the second-index
lifts above \(0,u\) in orbits of size \(2^j\).  Determinant primitivity
makes cyclotomic Galois act freely on the resulting orbits with order
\(2^j\).  The odd-order roots contribute only the fixed unramified data.
There is no additional choice of simultaneous sign after fixing the
old tuple.  As a degree check, the \(2^{2j}\) cusp orbits, residue degrees
\(2^j\), and ramification degrees \(2^j\) account for all \(2^{4j}\)
division tuples.

Set \(x=a_1/M_j\), \(y=a_2/M_j\).  Use the width-order
\(M_j\operatorname{ord}_q\).  The residue of the untwisted cup is the
Kummer class of the tame symbol, hence of \(\zeta_{M_j}^{h}\), where
\[
                  h\equiv\pm uJ_d(y)a_1J_c(x)\pmod{2^j}.
\]
The coefficient signs have exponents from \(-1\) divisible by \(2^j\).
The inverse twist uses the root \(\zeta_{2^j}^{\pm a_1u}\).
Consequently the twisted residue on \(G_{K_j}\) is the additive character
\[
          \pm h(a_1u)^{-1}\frac{\chi_{\mathrm{cy}}-1}{M_j}
                       \pmod{2^j}.
\]
Here \(\chi_{\mathrm{cy}}\) is the dyadic cyclotomic character; the odd
part of the roots is already fixed.  Transfer to \(G_{K_0}\) replaces
the cyclotomic expression by
\[
 \frac{\chi_{\mathrm{cy}}^{\,2^j}-1}{M_j}
               \equiv \frac{\log\chi_{\mathrm{cy}}}{M}\pmod{2^j}.
\]
The equality on the cyclotomic quotient is the abelian transfer
formula, and the congruence follows from \(v_2(M)\geq2\).
Residue trace sums these residue-field transfers without a
ramification-index multiplier.  Bernoulli distribution gives
\[
 \sum_{a_1}J_c(a_1/M_j)=J_c(1/M),\qquad
 \sum_{a_2}J_d(a_2/M_j)=J_d(1/M).
\]
The limiting residue, also the residue of the edge class by localization,
is therefore
\[
                  \pm J_c(1/M)J_d(1/M)
                              \frac{\log\chi_{\mathrm{cy}}}{M}.
\]

Étale and de Rham residues for \(H^1(1)\) compare by sending the
logarithmic differential of a parameter to residue one.  The dual
exponential on \(H^1(K_0,\mathbb Q_2)\) sends
\(\log\chi_{\mathrm{cy}}\) to \(\pm1\).  This follows by pairing with
Kummer classes of exponential units and local reciprocity, with field
trace in the de Rham pairing.

On the Eisenstein side, the constant term of \(E_{[x\tau+y]}\), for
\(0<x<1\), is \(-B_1(x)\).  One can obtain it by integrating in the
second summation coordinate before the prefactor: the constant
coefficient is the value at \(w=0\) of
\[
 -i\frac{\sqrt\pi\,\Gamma(1/2+w)}{\Gamma(1+w)}
 (\operatorname{Im}\tau)^{-w}
       \bigl[\zeta(2w,x)-\zeta(2w,1-x)\bigr],
\]
where \(\zeta(s,x)\) is the Hurwitz zeta function.  The two brackets
in \eqref{eq:K2} thus have constant terms
\(-J_c(1/M)\) and \(-J_d(1/M)\), both nonzero for the initial choices.
Since \(dq/q\) has residue \(M\) at this cusp, the proportionality
factor is exactly \(\pm M^{-2}\).  The orientation is common to all
components.

\emph{Arbitrary smoothing integers.}
In multiplicative group notation, divisor and norm normalization give
\[
 (c_0^2-[c_0]^*)\,{}_c\theta_E
                  =(c^2-[c]^*)\,{}_{c_0}\theta_E.
\]
Multiplying the first label by \(c\) in \eqref{eq:K1} changes the inverse
twist basis by \(c^{-1}\).  The induced operator is therefore
\(c^2-c\,\mathrm{mult}_c^*\).  Apply the two smoothing operators and
cancel the nonzero scalar \(c_0^2-c_0\) for
\(c_0\equiv1\pmod M\); do the same on the other row.  This is a
cancellation in the rational differential identity.  It makes no
assertion of integral divisibility by that scalar.  Row multiplication
may permute components, which is why a common calibration was used.
The formula follows for every \(c,d\) prime to \(6M\).
\end{proof}

\subsection{The norm relations}

\begin{proposition}[Good-prime norms and rectangular filling]
\label{prop:siegel-norms}
For either order pattern, adjoining a fresh prime
\(\ell\nmid cdn_2L\) to both orders and to \(M\) gives, after the
elliptic projection, the rational norm multiplier
\[
                   1-\frac{a_\ell(f)[\ell]}{\ell}
                               +\frac{[\ell]^2}{\ell}.
\]
The trace defining this norm is integral.  At complete splitting in a
lower abelian field the scalar is
\(\det(1-\rho_T(\mathrm{Fr}_\ell))/\ell\).

In rectangular order, let \(q_0^a\parallel b\), \(q_0\nmid b_1\).
At a finite even field character \(\nu\) of conductor dividing \(M\),
filling \(q_0^a\) into the first row gives the rational identity
\[
 z_{\mathrm{filled}}
       =\bigl(1-a_{q_0}(f)\nu(q_0)/q_0\bigr)
                                      z_{\mathrm{before}}.
\]
\end{proposition}

\begin{proof}
\emph{The fresh good prime.}
Before requiring independence at \(\ell\), the two division sets are
Cartesian.  Their norm is the old cup \eqref{eq:K1}, with mark,
inverse twist, and lower root agreeing, also through the dyadic tower.
Subtract the lifts whose \(\ell\)-parts both lie in a fixed order-\(\ell\)
line \(I_0\), sum over these lines, and add back \(\ell\) times the
double-prime-to-\(\ell\) lift.  This is inclusion--exclusion: the zero
pair belongs to all \(\ell+1\) lines.

For a subtracted line put
\(\rho:E\to\widetilde E=E/I_0\).  Distribution identifies the summed
cup with that on \(\widetilde E\) at the unique prime-to-\(\ell\)
points \(\widetilde P,\widetilde Q\) mapping to the old points by
\(\rho^\vee\); the relevant lifts were their \(\rho\)-preimages.
Reversing the description gives old valid tuples on \(\widetilde E\),
with arbitrary line \(\ker\rho^\vee\), mapping to \(E\) by quotient and
pushing the mark forward.  Relative to the source, the Weil roots
downstairs are raised to \(\ell\).  This gives the field translation
\([\ell]\) and the inverse-twist scalar \(1/\ell\).

This is the good Hecke correspondence, and it acts on the marked
elliptic push by \(a_\ell(f)\).  Indeed, pull back the relative forms
\(\Pi^*\omega,\Pi^*\bar\omega\) along the quotient and sum on the source;
\(F,\bar F\) are the corresponding Hecke eigenforms.  Equality holds
also for integrals along the lifted paths, with zero endpoint-function
coordinates.  It is consequently an equality of rational relative
Betti maps, and then of rational étale maps.  Thus no additional
boundary projection is needed.

The added-back lifts use multiplication by \(\ell\), of degree
\(\ell^2\), and roots raised to \(\ell^2\).  A possible diamond operator
on the mark acts trivially on this push, by the same relative-form
calculation.  The resulting multiplier is the one in the statement.
All maps used before and after the edge map are ordinary pullbacks
and unnormalized traces on fiber products.  At complete splitting,
\([\ell]=1\), and
\[
 1-\frac{a_\ell(f)}{\ell}+\frac1\ell
                  =\frac{\det(1-\rho_T(\mathrm{Fr}_\ell))}{\ell}.
\]

When integral derivative comparisons require removal of an error that
vanishes rationally, one common nonzero integer suffices in the non-CM
case.  Indeed, over the abelian fields in question, the torsion of
\(H^1(T)\) has bounded exponent.  By the open Tate image
\cite{Serre72}, the commutator image contains a fixed deep special-linear
congruence subgroup; its invariants on
\((T\otimes\mathbb Q_2)/T\) have bounded exponent.  The usual
cohomology sequence gives the asserted bound on \(H^1(T)\)-torsion.

\emph{Filling a bad-prime row.}
Keep \(Q\) and take the \(q_0^a\)-division lifts of \(P\).  The mark,
the determinant map to the \(M\)-th roots, and the inverse-root twist
agree.  Without the rank condition at \(q_0\), their sum is the old cup.
A discarded lift is one whose \(q_0\)-part has its
\([q_0^{a-1}]\)-image in the order-\(q_0\) line \(I_Q\) of the
\(Q\)-subgroup.  Put \(\rho:E\to\widetilde E=E/I_Q\).
These discarded lifts are exactly the preimages under
\(\rho[q_0^{a-1}]\) of the unique prime-to-\(q_0\) point
\(\widetilde P\) satisfying \(\rho^\vee\widetilde P=P\).
Their norm is \({}_c g_{\widetilde P}\).

Express \({}_d g_Q\) by the Kummer sum over
\(\rho^\vee\widetilde Q=Q\).  Each solution has order \(n_2\).
Indeed, by counting,
\(\rho^\vee\widetilde E[q_0^a]\) is the subgroup of
\(E[q_0^a]\) whose \([q_0^{a-1}]\)-image belongs to \(I_Q\).
Primitivity of \(Q\) also gives
\(I_{\widetilde Q}\ne\ker\rho^\vee\).
These descriptions continue to hold at the upper dyadic stages in
\eqref{eq:K1}.

In reverse, the discarded sum is therefore indexed by old valid tuples
on \(\widetilde E\) and all lines \(H_0\ne I_{\widetilde Q}\), with
quotient by \(H_0\).  Because \(l=(Ns,1)\), the line
\(I_{\widetilde Q}\) is also the order-\(q_0\) line of the source mark.
The roots are raised to \(q_0\), which is a unit on the shared field
and twist indices.  This marked correspondence is \(U_{q_0}\):
at mark \(1/L\), its branches are
\(\tau\mapsto(\tau+i)/q_0\), \(i\bmod q_0\), omitting exactly the
line of the mark.  It acts on \(F\) by \(a_{q_0}(f)\), since \(A\)
is a power of 2.  The translation of the root and the inverse twist
thus give the stated factor \(a_{q_0}(f)\nu(q_0)/q_0\) for the
discarded sum.  Subtracting it proves the filling identity.
\end{proof}

\subsection{The unnormalized elliptic period law}

\begin{proposition}[The period formula]\label{prop:siegel-period}
Use square order or first-row-unfilled rectangular order \(b_1=1\).
Suppose \(c,d\equiv1\pmod L\), and let \(\nu\) be an even finite
character of \(\operatorname{Gal}(\mathbb Q(\mu_M)/\mathbb Q)\).
For the projected, weighted class \(z_\nu\), the coefficient of its
dual exponential relative to \(\omega\) is
\[
 \pm\Delta_\nu\,
       \frac{L^{(M)}(f,\nu,1)}{\Omega_0}\,\delta_\xi(a),
 \qquad
 \Delta_\nu=(c^2-c\nu(c)^{-1})(d^2-d\nu(d)^{-1}).
                                                        \tag{K4}\label{eq:K4}
\]
Here \(L^{(M)}(f,\nu,s)\) twists the Fourier coefficients by \(\nu\)
and omits every prime dividing \(M\).  Scalars are compared in the
algebraic differential realizations and the fixed embeddings.  The
sign is common across the symbols.  No character-orbit average occurs.
\end{proposition}

\begin{proof}
\emph{The square-order trace.}
By Proposition~\ref{prop:siegel-reciprocity}, after tracing to the marked
level the coefficient of \(dq/q\) is
\[
 \frac{\pm\Delta_\nu}{(2\pi i)^2}
 \sum_{\substack{
 U=\left(\begin{smallmatrix}h&j\\m&n\end{smallmatrix}\right)
       \in M_2(\mathbb Z),\;(\det U,M)=1\\
                         lU\equiv(0,1)\bmod L}}
       \frac{\nu(\det U)}{(h\tau+j)(m\tau+n)}.
                                                        \tag{K5}\label{eq:K5}
\]
Regularize by multiplying the summand by
\[
         y_\tau^{u+v}|h\tau+j|^{-2u}|m\tau+n|^{-2v},
         \qquad y_z=\operatorname{Im}z,
\]
and continuing to \((u,v)=(0,0)\).  The row labels in a fiber over
the mark \(1/L\) satisfy exactly these congruences, for every embedding
of the determinant root; the target mark stays the same as that
embedding varies.  The Weil pairing is the determinant up to the
common orientation, which has no effect for even \(\nu\).
There is no simultaneous minus identification within this marked
fiber.  The two rescaled lattice sums contribute \(M^2\), canceling
the \(M^{-2}\) in \eqref{eq:K3}.  Multiplication of a row label by
\(c\) or \(d\) preserves the mark and, after reindexing its determinant
weight, produces the corresponding factor of \(\Delta_\nu\).

\emph{Hecke unfolding.}
Pair the sum with \(\bar F\) in \(dx\,dy\).
For positive determinant \(b_0\), unfold by the right
\(\Gamma_1(L)\)-action.  The substitution \(z=U\tau\) contributes
\(b_0^{-1-u-v}\) times
\[
 \mathcal J_{F|_2U^{-1}}(u,v)=
 \int_{\mathcal H}\overline{(F|_2U^{-1})(z)}
               \frac{y_z^{u+v}}{z|z|^{2u}}\,dx_z\,dy_z,
\]
where the weight-two slash operator includes the determinant factor.
The matrices \(\operatorname{adj}(\xi U)\) give the left
\(\Gamma_1(L)\)-cosets of determinant \(b_0\) with first column
\((1,0)^{\mathrm{tr}}\bmod L\).  These map bijectively to the usual
\(\Gamma_0(L)\) Hecke cosets, where the first entry is allowed to be
any unit: because \((b_0,L)=1\), a left \(\Gamma_0(L)\)-matrix can
normalize that entry to 1, uniquely modulo \(\Gamma_1(L)\).
The slash sum is therefore \(a_{b_0}(f)F|_2\xi\).
There is no extra \(-I\) stabilizer in \(\Gamma_1(L)\).

For negative determinants, flip the first row.  This gives the
subtraction with
\[
                  \xi'=J\xi J,\qquad J=\operatorname{diag}(-1,1).
\]
Thus, without the prefactor in \eqref{eq:K5}, the paired expression is
\[
 L^{(M)}(f,\nu,1+u+v)
       \bigl(\mathcal J_{F|_2\xi}
                    -\mathcal J_{F|_2\xi'}\bigr)(u,v).
\]

\emph{Continuation and the archimedean factor.}
Start in the region \(v>u>1\).
Each Kronecker series has moderate-growth meromorphic continuation
and is regular at zero in this situation; pairing with a cusp form
therefore permits continuation.  In the unfolded expression, first
hold \(v>3\) fixed and continue to \(u=0\).
For \(G=F|_2\xi\) or \(F|_2\xi'\),
\[
                  |G(z)|\leq C/y_z,
\]
and \(G\) has zero periodic mean and decays at infinity.  The integral
converges for \(0<\operatorname{Re}u<v\).
On horizontal tails, integrate by parts using a bounded periodic
primitive, which is also \(O(1/y_z)\) at small height.  This permits
the limit \(u=0\).  On bounded horizontal intervals, the displayed
growth bound, multiplied by \(y_z^{u+v}\), gives local domination.

At \(u=0\), the integral in \(x_z\), apart from \(y_z^v\), is
\(-2\pi i\,\overline{G(2iy_z)}\): this follows termwise from the
positive Fourier modes.  Cuspidality then permits continuation of
the remaining integral to \(v=0\), yielding
\[
           \mathcal J_G(0,0)
                 =\pi\,\overline{\int_0^{i\infty}G(z)\,dz}.
\]
Write \(I=\int_0^{i\infty}(F|_2\xi)(z)\,dz\).
The integral for \(\xi'\) is \(-\bar I\), since the coefficients
of \(f\) are real.  Since
\[
        \frac{dq}{q}\wedge\frac{d\bar q}{\bar q}
                                   =-8\pi^2i\,dx\,dy,
\]
the wedge pairing of the traced differential with \(\Pi^*\bar\omega\)
is
\[
       \pm2i\,\operatorname{Im}(A_\Pi)\,
                     \Delta_\nu L^{(M)}(f,\nu,1),
       \qquad A_\Pi=c_F(2\pi i)I.
\]

\emph{Projection to the elliptic differential.}
The relative pullback has zero cusp-function coordinates.
Cutoff primitives vanishing at the cusps introduce no boundary term
when paired with a differential having at worst logarithmic poles.
The push is dual to this pullback, with the same trace convention on
both curves.  Write
\(\delta_\xi=s_0\omega+\bar s_0\bar\omega\).  Then
\[
 \int_{E'}\omega\wedge\delta_\xi=A_\Pi,\qquad
 2i\,\operatorname{Im}(A_\Pi)
   =(s_0+\bar s_0)\int_{E'}\omega\wedge\bar\omega,\qquad
 s_0+\bar s_0=\frac{\delta_\xi(a)}{\Omega_0}.
\]
The projected differential is on the Hodge line.  These identities
prove \eqref{eq:K4} in square order, with consistent sign for all
symbols.

This comparison of complex and dyadic scalars is algebraic:
the modular differentials with the specified weights are algebraic,
and trace and push are algebraic de Rham operations.  After restriction
to a field trivializing the character, dual exponential uses the same
weighted sum.  Its local duality uses field trace, not the average of
that trace.

\emph{Rectangular order.}
Apply the filling identity of Proposition~\ref{prop:siegel-norms}
one full prime power \(q_0^a\parallel b\) at a time.
This fills the first row to square level \(bM\), with determinant
weight inflated to that level.  Each factor
\[
                         1-a_{q_0}(f)\nu(q_0)/q_0
\]
is nonzero at a finite character, since \(|a_{q_0}(f)|\leq1\).
It is exactly the extra omitted bad-prime Euler factor at \(q_0\).
Canceling these factors in the rational period identity for square
level gives \eqref{eq:K4} with omissions at \(M\), as asserted.
This cancellation proves a rational differential formula.  The
integral fresh-prime trace for the rectangular system was constructed
separately in Proposition~\ref{prop:siegel-norms}.
\end{proof}

\subsection{Choices of symbols and positive quadratic twists}

\begin{lemma}[A nonzero rectangular symbol]\label{lem:siegel-rectangular-symbol}
In rectangular order one can choose \(A=2^k\) and
\(\xi\in\mathrm{SL}_2(\mathbb Z)\) with lower row \((Ns,1)\) so that
\(\delta_\xi(a)\ne0\).
\end{lemma}

\begin{proof}
Take the upper row to be \((1,0)\).  After the degeneracy, the path
runs from \(0\) to \(2^k/(Ns)\).
The Fricke relation for \(f\) identifies the integral of \(2\pi i f\)
along it, up to sign, with the integral from \(\infty\) to
\(-s/2^k\).  For some \(k\), these latter integrals cannot all be real.
Indeed, Rohrlich's finite-support twist nonvanishing
\cite{Rohrlich88}, in the form \cite[Theorem~13.5(2)]{Kato04}, applies
with prime support \(\{2\}\).  One can equivalently fix the odd
quadratic character \(\eta=\chi_{-4}\) and vary primitive even
characters of the real cyclotomic tower: multiplying a character of
conductor \(2^k\), \(k\geq3\), by \(\eta\) preserves that conductor
and makes it odd.  Thus there is a primitive odd character \(\nu\)
of sufficiently high dyadic conductor with
\(L(f,\bar\nu,1)\ne0\).

Set \(\lambda(r)=\int_\infty^r2\pi i f(z)\,dz\).  The primitive
additive Gauss--Mellin formula is
\[
 \sum_{s\bmod 2^k}\nu(s)\lambda(-s/2^k)
       =\nu(-1)\tau(\nu)L(f,\bar\nu,1),\qquad
 \tau(\nu)=\sum_{s\bmod2^k}\nu(s)e^{2\pi i s/2^k}.
\]
Its right-hand side is nonzero.  If all the indicated integrals were
real, the identity \(\lambda(-r)=\overline{\lambda(r)}\), which
follows from the real Fourier coefficients, would make them unchanged
under \(s\mapsto-s\).  Their sum against the odd character \(\nu\)
would then vanish.  This contradiction gives a nonzero imaginary part.
The path pairing in the proof of
Proposition~\ref{prop:siegel-period} then gives
\(\delta_\xi(a)\ne0\).
\end{proof}

\begin{corollary}[Primitive evaluation and positive twisting]
\label{cor:siegel-positive-twist}
For the square construction with \(A=1\), one may use an integral
combination of symbols with \(\delta(a)=1\) and common smoothing
integers.  If a fixed positive fundamental quadratic character
\(\chi_h\), prime to \(2N\), is included in the field weight, the
integral coefficient identification with the twist \(E'^{\,h}\)
transports \eqref{eq:K4} to the corresponding twist-period formula,
with its actual smoothing factors and omissions.
\end{corollary}

\begin{proof}
The first assertion is Lemma~\ref{lem:siegel-lattice} and linearity.
Include the conductor of \(\chi_h\) in \(M\).
The weighted coefficient with this fixed factor identifies integrally
with the coefficient for \(E'^{\,h}\) by the geometric twist
isomorphism.  Choose this isomorphism compatibly over \(\mathbb R\);
positivity of \(h\) makes the primitive real-invariant cycles correspond
up to sign.  If \(\omega\) pulls back to \(b_h\omega_h\), then the
transported exponential coefficient relative to \(\omega_h\) is
multiplied by \(b_h\), while
\[
                         \Omega_{0,h}=\Omega_0/|b_h|.
\]
The omitted-support \(L\)-value on the twist uses the remaining
character.  Thus both sides of \eqref{eq:K4} transport with the same
factor, up to the common sign.  The coefficient construction is the
weighted integral trace, so no division by a quadratic projector is
made.
\end{proof}

The period in \eqref{eq:K4} remains \(\Omega_0\).
The factor \(\#\pi_0(E'(\mathbb R))\) converting it to the total real
period enters through real Kummer cohomology in the arithmetic
determinant comparison.
